% ===== 第八章 附录：求解代码（约4页） =====
\section*{八、附录：求解代码与数值实现}

本节给出全文数值求解的核心代码（Python），对应仓库 $engines/cumcm2024a\_solver.py$ 与 $engines/run\_all\_2024A.py$。全部结果、参数与算法细节见 $02\_execution\_log.json$。

\subsection*{8.1 求解核心库（弧长闭式、逆弧长、碰撞检测）}

\begin{lstlisting}[language=Python, caption={cumcm2024a\_solver.py（节选）}]
import numpy as np, json, os
from scipy.integrate import quad

HEAD_H2H=2.86; BODY_H2H=1.65; N=224
GAPS=np.array([HEAD_H2H]+[BODY_H2H]*222)      # 段 p--p+1 板长
CHAIN=GAPS.sum()                              # 369.16

def spiral_pt(theta,a):
    r=a*theta; return np.stack([r*np.cos(theta), r*np.sin(theta)],axis=-1)

def s_arc(th,a):                                # 弧长闭式 (精确)
    t=np.asarray(th,dtype=float)
    return a/2*(t*np.sqrt(1+t*t)+np.arcsinh(t))

def inv_arc(s,a,hi=6000.0):                     # 500步二分逆弧长
    lo=1e-14
    for _ in range(500):
        m=0.5*(lo+hi)
        if s_arc(m,a)<s: lo=m
        else: hi=m
    return 0.5*(lo+hi)

def arc_positions(s_head, a, D):
    th=np.array([inv_arc(max(s_head-D[k],1e-9),a) for k in range(N)])
    return spiral_pt(th,a)

def arc_pos_vel(s_head, a, D, ds=1e-7):
    P=arc_positions(s_head,a,D)
    P2=arc_positions(s_head+ds,a,D)
    dr=(P2-P)/ds
    v=np.linalg.norm(dr,axis=-1)
    return P, v, dr

def seg_intersect(aaa,bbb,ccc,ddd,tol=1e-9):
    def cross(o,p,q): return (p[0]-o[0])*(q[1]-o[1])-(p[1]-o[1])*(q[0]-o[0])
    d1=cross(ccc,ddd,aaa); d2=cross(ccc,ddd,bbb)
    d3=cross(aaa,bbb,ccc); d4=cross(aaa,bbb,ddd)
    return ((d1>tol and d2<-tol)or(d1<-tol and d2>tol))and\
           ((d3>tol and d4<-tol)or(d3<-tol and d4>tol))

def collision_segments(P, skip=1):
    for i in range(N-1):
        A,B=P[i],P[i+1]
        for j in range(i+2,N-1):
            C,DCS=P[j],P[j+1]
            if seg_intersect(A,B,C,DCS): return (i,j)
    return None
\end{lstlisting}

\subsection*{8.2 主求解流程（全量求解与交付）}

\begin{lstlisting}[language=Python, caption={run\_all\_2024A.py（主流程节选）}]
import numpy as np, json, os, pandas as pd
sys.path.insert(0, os.path.dirname(os.path.abspath(__file__)))
import cumcm2024a_solver as C

D=np.array([0.0]+[2.86+(k-1)*1.65 for k in range(1,224)])  # D[223]=369.16
CHAIN=D[-1]
A1=0.55/(2*np.pi); S0=C.s_arc(32*np.pi,A1)                 # S0=442.59026

# Q1: 0..300s 逐秒
rec1=[]
for t in range(0,301):
    s=S0-t; P,v,dr=C.arc_pos_vel(s,A1,D)
    for h in range(C.N):
        rec1.append({"t":t,"handle":h,"x":round(float(P[h,0]),6),
                     "y":round(float(P[h,1]),6),"v":round(float(v[h]),6)})
df=pd.DataFrame(rec1); df.to_excel(SUBD+"/result1.xlsx",index=False)

# Q2: 二分求 t* (双判据)
def collide(t):
    s=S0-t
    if s<CHAIN: return "tail_past_center"
    P=C.arc_positions(s,A1,D)
    return ("seg_"+str(k)) if (k:=C.collision_segments(P)) else None
lo,hi=0.0,S0-CHAIN
for _ in range(70):
    md=0.5*(lo+hi)
    if collide(md) is None: lo=md
    else: hi=md
t2=hi                                   # t*=73.4303s

# Q4: S形调头两圆弧 R1=2R2 由相切约束数值解出
R2,b = solve_R2(A4,TH_MAR)              # R2=1.5027, R1=3.0054, b=3.0215
path,E,Xend,n_in,turn = q4_composite_path(R2,b)
# 逐秒 -100..100s 布置整链 224 把手并写 result4.xlsx
\end{lstlisting}

\subsection*{8.4 问题四复合路径构造（S形调头）}
\begin{lstlisting}[language=Python, caption={Q4 复合路径解析构造（节选）}]
def q4_composite_path(R2, b, n_in=1500, n_arc=120):
    R1=2*R2; thE=4.5/A4
    E=spiral_pt(thE, A4); X=-E                     # 入口E/出口X 均在 r=4.5 边界(中心对称)
    din=dP(thE,A4); v_out=-din/np.linalg.norm(din); n=[-v_out[1],v_out[0]]
    th_in=np.linspace(TH_MAR, thE, n_in)           # 盘入螺线: 外(θ大)→内(θ小), 末端=E
    P_in=spiral_pt(th_in, A4)
    C1=E+R1*(-n)                                   # 前段圆弧(半径R1)与盘入螺线在E相切
    a0=np.arctan2((E-C1)[1],(E-C1)[0])
    arc1=[C1+R1*[np.cos(a0-k*b/(n_arc-1)),np.sin(a0-k*b/(n_arc-1))] for k in range(n_arc)]
    J=arc1[-1]
    vJ=rot(v_out,-b); nvJ=[-vJ[1],vJ[0]]
    C2=J+R2*nvJ                                    # 后段圆弧(半径R2)与前一圆弧外切
    aJ=np.arctan2((J-C2)[1],(J-C2)[0])
    arc2=[C2+R2*[np.cos(aJ+k*b/(n_arc-1)),np.sin(aJ+k*b/(n_arc-1))] for k in range(n_arc)]
    P_out=-spiral_pt(np.linspace(thE,TH_MAR,n_in), A4)   # 中心对称盘出
    return np.vstack([P_in,arc1,arc2,P_out]), E, X
# 相切约束数值解: R1=2R2, 前段切E, 后段切X => R2=1.5027, b=3.0215
\end{lstlisting}

\subsection*{8.5 问题四几何与间距表 [log:q4]}
表 \ref{tab:q4geo} 汇总问题四的关键几何量与间距指标。

\begin{table}[H]\centering
\caption{问题四关键量与间距验证[log:q4]}
\label{tab:q4geo}
\begin{tabular}{lcc}
\toprule
量 & 值 & 说明 \\
\midrule
盘入螺距 $p$ & 1.7 m & $a=p/2\pi$ \\
前段圆弧半径 $R_1$ & 3.0054 m & $R_1=2R_2$ \\
后段圆弧半径 $R_2$ & 1.5027 m & 相切约束解 \\
每弧张角 $\Delta$ & 3.0215 rad & 约 173.1° \\
调头弧长 $L_{\rm turn}$ & 13.621 m & $(R_1+R_2)\Delta$ \\
全程最小相邻间距 & 1.5683 m & 高于 1.50m 阈值 \\
有无碰撞 & 无 & $-100\sim100$s 扫描 \\
调头圆弧最大半径 & 4.500 m & 恰在 r≤4.5 空间内 \\
\bottomrule
\end{tabular}
\end{table}

\subsection*{8.6 完整性}
全部数值以 6 位小数写入 $result1/2/4.xlsx$，并同步写入 $02\_execution\_log.json$ 中的键（$q1.paper\_sample$、$q2.t\_star$、$q3.min\_pitch$、$q4.R1/R2/turn\_arc\_len$、$q5.beta\_max$），保证每个进入论文的数字可溯源。[log:q1.config, q2.config, q4.config]